\chapter{Complete v0.0.6 Rule Catalogue}
\label{chap:rules}

This chapter is generated from the conceptual structure of
\code{RULE_CATALOG} in the v0.0.6 source. It documents all 76 active concepts.
The catalogue has 12 categories; the frozen snapshot result table has 11
populated categories because \emph{legal} contains a review-range rule and no
legal occurrence was present in the source-snapshot aggregation. A catalogue
category must therefore not be confused with a non-zero result category.

For each rule, the table states the stable identifier, default severity, and
operational meaning. Exact thresholds remain configurable; the descriptions
below explain the predicate rather than claiming that every match is a defect.

\section{Security Rules (8)}
\begin{longtable}{p{4.2cm}p{1.4cm}p{8cm}}
\toprule Rule identifier & Level & Detection intent \\ \midrule\endhead
\code{become-scope-too-broad} & warning & Detects play-level privilege escalation that applies root scope to every task. \\
\code{credentials-leaking-in-log} & error & Detects a task handling sensitive data without known \code{no_log: true} protection. \\
\code{default-credentials} & error & Detects credential fields assigned known default or placeholder values. \\
\code{hardcoded-credentials} & error & Canonical concept receiving secret-pattern evidence from GitLeaks. \\
\code{insecure-permissions} & error & Detects managed resources whose mode grants unsafe access to other users. \\
\code{missing-validate-certs} & warning & Detects disabled TLS verification for HTTPS-capable modules. \\
\code{plaintext-sensitive-variable} & warning & Detects sensitive variable names assigned plaintext literals. \\
\code{vault-inline-overload} & note & Detects many inline vault values where a dedicated vault file may be clearer. \\
\bottomrule
\end{longtable}

\section{Correctness Rules (14)}
\begin{longtable}{p{4.2cm}p{1.4cm}p{8cm}}
\toprule Rule identifier & Level & Detection intent \\ \midrule\endhead
\code{broken-role-dependency} & error & Detects missing, cyclic, or self-referential local role dependencies. \\
\code{invalid-numeric-value} & error & Detects supported numeric module values outside their valid range. \\
\code{invalid-task-attribute-type} & error & Detects common task attributes represented with the wrong YAML type. \\
\code{missing-bool-filters} & warning & Detects conditions comparing values with quoted boolean strings. \\
\code{missing-handler} & error & Detects \code{notify} targets without a resolvable handler definition. \\
\code{non-idempotent-command} & warning & Detects mutating command/shell tasks without a modelled idempotency guard. \\
\code{overriden-keys} & error & Detects duplicate YAML mapping keys that silently discard an earlier value. \\
\code{overriding-facts} & warning & Detects a \code{set_fact} name reassigned with a different expression. \\
\code{role-variable-precedence-collision} & warning & Detects the same role variable in both defaults and vars. \\
\code{shell-environment-export} & warning & Detects shell \code{export} where Ansible's environment mapping should be used. \\
\code{unguarded-variable} & warning & Classifies a Jinja use from visible definition, guard, runtime source, order, and uncertainty evidence. \\
\code{unscoped-shell-variable} & warning & Detects shell-variable use not declared through the task environment. \\
\code{variable-use-before-definition} & warning & Detects a use before a later registered result or persistent fact producer. \\
\code{yaml-truthy-value} & warning & Detects non-canonical YAML truth values such as \code{yes}/\code{no}. \\
\bottomrule
\end{longtable}

\section{Maintainability Rules (15)}
\begin{longtable}{p{4.2cm}p{1.4cm}p{8cm}}
\toprule Rule identifier & Level & Detection intent \\ \midrule\endhead
\code{complex-expression} & warning & Detects Jinja or conditional expressions above configured complexity. \\
\code{complex-playbook} & warning & Detects aggregate playbook complexity above threshold. \\
\code{deprecated-modules} & warning & Detects deprecated module use through the compatibility map. \\
\code{duplicated-code} & warning & Detects structurally repeated task blocks. \\
\code{excessive-nesting} & warning & Detects deeply nested block/rescue/always structures. \\
\code{high-coupling} & warning & Detects unusually many role/import relationships. \\
\code{incomplete-task} & warning & Detects unresolved TODO/FIXME markers in configuration code. \\
\code{insufficient-modularization} & warning & Detects large multi-play files that should be split. \\
\code{large-play} & warning & Detects plays exceeding size thresholds. \\
\code{large-playbook} & warning & Detects playbooks exceeding aggregate size thresholds. \\
\code{large-task} & warning & Detects task mappings exceeding size thresholds. \\
\code{redundant-default-value} & note & Detects an assignment repeating a role-default value. \\
\code{shell-instead-of-module} & warning & Detects shell/command forms replaceable by modelled native modules. \\
\code{unused-handler} & warning & Detects defined handlers that are never notified. \\
\code{unused-registered-variable} & note & Detects registered results without a visible consumer. \\
\bottomrule
\end{longtable}

\section{Style Rules (13)}
\begin{longtable}{p{4.2cm}p{1.4cm}p{8cm}}
\toprule Rule identifier & Level & Detection intent \\ \midrule\endhead
\code{duplicate-task-names} & note & Detects repeated task names within the applicable scope. \\
\code{inconsistent-naming-convention} & note & Detects names dominated by special characters or inconsistent forms. \\
\code{long-statement} & note & Detects YAML source lines above the configured limit. \\
\code{long-task-name} & note & Detects task names above the configured limit. \\
\code{missing-file-extension} & note & Detects file references without an extension where one is expected. \\
\code{missing-tags} & note & Optional policy detecting tasks without explicit or inherited tags. \\
\code{unnamed-play} & note & Detects plays without \code{name}. \\
\code{unnamed-task} & note & Detects tasks without \code{name}. \\
\code{unnecessary-abstraction} & note & Detects a block wrapping only one task. \\
\code{unnecessary-set-fact} & note & Detects a fact assignment with no modelled consumer. \\
\code{unused-include-vars} & note & Detects included variable data with no modelled use. \\
\code{variable-naming-convention} & note & Detects role variables without the role-name prefix. \\
\code{yaml-formatting} & note & Detects indentation, whitespace, and line-ending problems. \\
\bottomrule
\end{longtable}

\section{Compatibility Rules (4)}
\begin{longtable}{p{4.2cm}p{1.4cm}p{8cm}}
\toprule Rule identifier & Level & Detection intent \\ \midrule\endhead
\code{deprecated-module-option} & warning & Detects deprecated options for recognised modules. \\
\code{fqcn-migration} & note & Detects short module names with a known fully qualified form. \\
\code{loop-legacy-syntax} & warning & Detects legacy \code{with_*} iteration instead of \code{loop}. \\
\code{redirected-module} & note & Detects legacy module names redirected to a collection FQCN. \\
\bottomrule
\end{longtable}

\section{Documentation Rules (6)}
\begin{longtable}{p{4.2cm}p{1.4cm}p{8cm}}
\toprule Rule identifier & Level & Detection intent \\ \midrule\endhead
\code{commit-message-format} & note & Checks the latest commit message against review conventions. \\
\code{galaxy-meta-incomplete} & note & Detects missing required Galaxy metadata fields. \\
\code{missing-bug-reference} & note & Detects a bug-fix commit without a bug reference. \\
\code{missing-release-note} & note & Detects a configuration change without an associated release note. \\
\code{role-readme-absent} & note & Detects roles without a README. \\
\code{undocumented-complex-task} & note & Detects complex tasks without a nearby explanatory comment. \\
\bottomrule
\end{longtable}

\section{Structure Rules (4)}
\begin{longtable}{p{4.2cm}p{1.4cm}p{8cm}}
\toprule Rule identifier & Level & Detection intent \\ \midrule\endhead
\code{empty-abstraction} & warning & Detects empty task files or nested task abstractions. \\
\code{missing-abstraction} & warning & Detects large task collections not split into roles or includes. \\
\code{unexpected-task-location} & warning & Detects task/handler documents in incompatible role directories. \\
\code{unstructured-module} & warning & Detects roles that do not follow the conventional directory layout. \\
\bottomrule
\end{longtable}

\section{Reliability Rules (3)}
\begin{longtable}{p{4.2cm}p{1.4cm}p{8cm}}
\toprule Rule identifier & Level & Detection intent \\ \midrule\endhead
\code{read-only-command-change} & note & Detects read-only commands that may report changed on every run. \\
\code{silent-error-suppression} & warning & Detects ignored failure without retaining a result for later control. \\
\code{unsafe-file-write} & warning & Detects file operations disabling atomic-write protection. \\
\bottomrule
\end{longtable}

\section{Testing Rules (3)}
\begin{longtable}{p{4.2cm}p{1.4cm}p{8cm}}
\toprule Rule identifier & Level & Detection intent \\ \midrule\endhead
\code{missing-test-update} & warning & Detects configuration changes without corresponding test changes. \\
\code{molecule-test-absent} & note & Detects a role without a Molecule scenario. \\
\code{untested-override-change} & warning & Detects defaults/vars changes without an override-test change. \\
\bottomrule
\end{longtable}

\section{Dependencies Rules (2)}
\begin{longtable}{p{4.2cm}p{1.4cm}p{8cm}}
\toprule Rule identifier & Level & Detection intent \\ \midrule\endhead
\code{invalid-depends-on-reference} & warning & Detects malformed \code{Depends-On} review trailers. \\
\code{unpinned-ansible-dependency} & warning & Detects role/collection dependencies without a fixed version. \\
\bottomrule
\end{longtable}

\section{Cohesion Rules (2)}
\begin{longtable}{p{4.2cm}p{1.4cm}p{8cm}}
\toprule Rule identifier & Level & Detection intent \\ \midrule\endhead
\code{incohesive-play} & note & Uses module-domain entropy to identify plays spanning unrelated concerns. \\
\code{incohesive-playbook} & note & Applies the same domain-diversity model at playbook scope. \\
\bottomrule
\end{longtable}

\section{Lint and Legal Rules (2)}
\begin{longtable}{p{4.2cm}p{1.4cm}p{8cm}}
\toprule Rule identifier & Level & Detection intent \\ \midrule\endhead
\code{ansible-lint-issue} & warning & Envelope for native ansible-lint evidence; the native rule remains in provider evidence. \\
\code{removed-copyright-notice} & warning & Review-range check detecting removal of an existing copyright notice. \\
\bottomrule
\end{longtable}

\section{How to Interpret a Rule Match}
A rule match has four independent dimensions. The \emph{concept} states what
was detected; \emph{severity} states default policy impact; \emph{semantic
status} states evidence confidence or runtime limitation; and \emph{provider
evidence} records who observed it. For example, a warning-level
\code{unguarded-variable} can be \code{UNCERTAIN} because a dynamic inventory
may define it. Severity must not be used as a substitute for semantic status,
and neither is a human true/false-positive label.
